将转移写成 $f_i=A_i+\max_{j<i}(k_jx_i+b_j)$，每个已求出的状态 $j$ 对应一条直线。先在 $x_i$ 处查询，再插入状态 $i$ 对应的直线。

斜率非递减时，使用单调斜率版的 \texttt{query}；查询坐标也非递减时，可调用 \texttt{mono}。斜率任意时使用任意斜率版。各次查询可以与插入交替进行，不需要预设坐标范围。若直线只在某个区间内有效，使用李超树的线段版本。
